% --- BEGIN MERGED BLOCK: methods appendix (method-appendix-semantic-ising.tex) ---
\section{Semantic and Ising Construction of Considerations}

This appendix documents the empirical construction of the consideration pool
used in the Seattle minimum-wage simulations. The goal of the procedure is not
to recover a complete theory of the original Polis discussion, but to convert a
large, irregular public-comment dataset into a compact set of considerations
that can be used in a controlled agent-based simulation. The design problem is
therefore one of disciplined reduction: preserve the main argumentative
contrasts in the discussion, avoid letting near-duplicate statements dominate
the model, and retain empirical information about which considerations were
endorsed together or rejected together.

\subsection{Data Source and Preprocessing}

The source case is the 2014 Seattle Polis conversation on the proposal to raise
the minimum wage to \$15 per hour. The open-data release contains 339
participants and 69 submitted statements. We restrict the statement pool to the
30 statements that passed moderation in the Polis dataset. Votes are coded as
$+1$ for agree, $-1$ for disagree, and $0$ for pass; missing votes are treated
as absent rather than as passes. Because support and opposition relations are
estimated from directional voting behavior, pass and missing responses are
excluded from pairwise co-vote calculations.

A statement must receive at least 34 directional votes, approximately ten
percent of the participant pool, to be eligible for selection. This threshold
removes one moderated statement with too little directional information for
stable pairwise association estimates. The remaining 29 statements form the
candidate set for semantic clustering and representative selection.

\subsection{Co-Vote Association}

For every pair of eligible statements $i$ and $j$, we estimate the association
between their vote profiles using Yule's $Q$ for two-by-two association tables
\citep{yule1912}. Let $A_{ij}$ be the number of participants who voted in the
same direction on both statements, either agree-agree or disagree-disagree, and
let $D_{ij}$ be the number who voted in opposite directions. The pairwise
association is
\[
Q_{ij} = \frac{A_{ij} - D_{ij}}{A_{ij} + D_{ij}},
\]
computed only when at least 10 participants cast directional votes on both
statements. $Q_{ij}=1$ indicates perfect co-endorsement or co-rejection,
$Q_{ij}=-1$ indicates perfect opposition, and $Q_{ij}=0$ indicates no net
directional association among participants who voted on both statements.

This statistic gives the model an empirical measure of how statements relate
within the actual participant response structure. It also avoids treating raw
agreement rates as if they were argument relations. A popular statement is not
therefore automatically a strong support, nor is an unpopular statement
automatically an attack. Support and attack relations depend on how statements
pattern together across participant votes.

On the eligible 29-statement set, the resulting off-diagonal $Q$ values have
mean $+0.099$, standard deviation $0.339$, and range $[-0.714, +0.803]$.
Approximately 59.4\% of pairwise associations are positive and 37.2\% are
negative, so the empirical graph contains substantial opposition structure
rather than just diffuse agreement.

\subsection{Semantic Clustering and Representative Selection}

The eligible statements are embedded with the \texttt{e5-large-v2} text
embedding model, part of the E5 family of weakly supervised contrastive text
embeddings \citep{wang2022e5}. We cluster the 29 statement embeddings using
hierarchical agglomerative clustering with average linkage and cosine distance,
fixing the number of clusters at $K=10$. The number of clusters is chosen to
produce a tractable consideration pool for the simulations while preserving
distinct argumentative themes in the original discussion: living wage, inflation
and prices, income inequality, small-business pressure, automation, guaranteed
income, service redesign, and systemic critique.

As a validation check, we compared archived factor-analytic statement
assignments to embedding-only clusters. The adjusted Rand index remains low
across models and cluster counts: at $K=10$, ARI is $0.098$ for
\texttt{nomic-embed-text}, $0.092$ for raw \texttt{e5-large-v2}, and $0.121$
for contextual \texttt{e5-large-v2}; across all tested specifications the ARI
range is $[-0.003, 0.138]$. At the same time, within-factor semantic coherence
is high: the mean within-factor cosine similarity is $0.935$, the weakest
factor mean is $0.904$, and all 10 archived factors exceed $0.85$. This is the
empirical reason for using embeddings to ensure topical coverage while letting
vote behavior, rather than embedding geometry alone, determine the association
structure.

Within each semantic cluster, we select one representative statement. The
representative is not chosen by agreement rate. Instead, for each statement $i$
in cluster $C(i)$, we compute its separation power:
\[
s_i = \max_{C' \neq C(i)}
\frac{1}{|C'|}\sum_{j \in C'} |Q_{ij}|.
\]
This score identifies the statement in a semantic cluster that most clearly
separates that cluster from at least one other cluster in the participant
co-vote structure. Ties are broken by the number of directional votes. This
rule is deliberately conservative: it keeps one semantically distinct
representative per cluster, and it chooses the representative whose vote
pattern contributes most to the empirical structure of disagreement.

The direction of each selected consideration is assigned by substantive
content, with $+1$ indicating a consideration that supports the minimum-wage
increase and $-1$ indicating a consideration that opposes or criticizes it.
This content coding is necessary because net agreement is not a reliable
measure of policy direction. Participants can agree with a critical statement
while supporting the policy, or agree with a systemic critique that does not
map cleanly onto simple support for the proposal.

\begingroup
\scriptsize
\setlength{\tabcolsep}{2pt}
\begin{longtable}{p{0.34in}p{0.34in}p{0.30in}p{0.32in}p{0.36in}p{0.40in}p{0.40in}p{3.35in}}
\caption{Selected considerations used in the simulation scenario.}\label{tab:selected-considerations}\\
\toprule
ID & Polis & Dir. & Clus. & Sep. & Cross & Pers. & Statement \\
\midrule
\endfirsthead
\toprule
ID & Polis & Dir. & Clus. & Sep. & Cross & Pers. & Statement \\
\midrule
\endhead
\bottomrule
\endfoot
C\_01 & 48 & Pro & 1 & .742 & .471 & .817 & Businesses, small or large, shouldn't be able to pay workers anything below a livable wage. Why should we subsidize businesses that depend on inadequate wages? If a business's product is good, then this shouldn't be a challenge. \\
C\_02 & 11 & Pro & 2 & .571 & .541 & .900 & Given inflation, unless workers fight they are continually being pushed down. I think we'll need to revisit these laws again and again to make sure the economy is working for poor families. \\
C\_03 & 32 & Con & 3 & .571 & .032 & .300 & \$15 min wage will lead to inflation. Higher income population is penalized for this reason and therefore I don't think it is fair. \\
C\_04 & 8 & Pro & 4 & .690 & .314 & .633 & During the campaign, I found out that 75\% of minimum-wage workers are 20 or older. I think it's the right thing to do to bring higher incomes to this group of people. \\
C\_05 & 45 & Pro & 5 & .739 & .400 & .734 & Something needs to be done to address income inequality and this is a good if imperfect start. \\
C\_06 & 46 & Con & 6 & .657 & .130 & .416 & Higher wages makes it super hard for small businesses to survive. In effect, mom and pop coffee shops with disappear. \\
C\_07 & 36 & Con & 7 & .458 & .429 & .768 & It's just going to speed up the adoption of robotics in industries with unskilled or low-skilled workers. \\
C\_08 & 39 & Pro & 8 & .431 & .500 & .852 & One step towards a world where a basic minimum income will be the way developed nations must function as the need for human workers collapses vs populace that ``needs a job.'' \\
C\_09 & 0 & Con & 9 & .234 & .414 & .750 & I imagine new businesses opening will run counter operations where guests grab their own silverware and bus their own tables. \\
C\_10 & 20 & Con & 10 & .556 & .061 & .333 & Reactionary liberal trash. Doesn't solve the problem. It's a systemic problem and this solution is like a band aid on a broken leg. \\
\end{longtable}
\endgroup

\subsection{Support and Attack Edges}

The support and attack graph is constructed from the same Yule's $Q$ matrix.
For each selected pair of considerations, we add an edge when $|Q_{ij}| \geq
0.2$ and at least 10 participants cast directional votes on both statements.
Negative associations are treated as attacks, because participants who endorse
one consideration tend to reject the other. Positive associations are treated
as supports, because participants tend to endorse or reject the two
considerations together. Edge strength is scaled as
\[
\omega_{ij} = 0.2 + 0.6\min(|Q_{ij}|,1).
\]

Attack edges are stored symmetrically in the scenario because empirical
opposition between two considerations is undirected: if consideration $i$
contradicts the vote profile of consideration $j$, then $j$ also contradicts
$i$. Support edges are retained as directed edges in the scenario object, but
their empirical source is the same symmetric co-vote association. The directed
storage convention allows the simulation engine to treat support relations as
an ordered input when needed, while the method itself should be read as a
co-endorsement measure rather than a causal claim.

\begingroup
\scriptsize
\setlength{\tabcolsep}{2pt}
\begin{longtable}{p{0.54in}p{0.45in}p{0.45in}p{0.36in}p{4.20in}}
\caption{Empirical support and attack relations among selected considerations.}\label{tab:consideration-edges}\\
\toprule
Relation & From & To & Str. & Interpretation \\
\midrule
\endfirsthead
\toprule
Relation & From & To & Str. & Interpretation \\
\midrule
\endhead
\bottomrule
\endfoot
Attack & C\_01 & C\_03 & .457 & Living-wage/subsidy argument versus inflation and fairness concern. \\
Attack & C\_01 & C\_06 & .320 & Living-wage/subsidy argument versus small-business survival concern. \\
Attack & C\_01 & C\_10 & .615 & Living-wage/subsidy argument versus systemic-critique rejection of the reform. \\
Attack & C\_02 & C\_03 & .543 & Inflation-protection argument for workers versus inflation and fairness concern. \\
Attack & C\_02 & C\_06 & .490 & Worker-protection argument versus small-business survival concern. \\
Attack & C\_02 & C\_10 & .512 & Worker-protection argument versus systemic-critique rejection of the reform. \\
Attack & C\_03 & C\_04 & .513 & Inflation/fairness concern versus older-worker income argument. \\
Attack & C\_03 & C\_05 & .534 & Inflation/fairness concern versus inequality-reduction argument. \\
Attack & C\_03 & C\_08 & .416 & Inflation/fairness concern versus basic-income transition argument. \\
Attack & C\_04 & C\_06 & .541 & Older-worker income argument versus small-business survival concern. \\
Attack & C\_04 & C\_07 & .384 & Older-worker income argument versus automation concern. \\
Attack & C\_04 & C\_10 & .494 & Older-worker income argument versus systemic-critique rejection of the reform. \\
Attack & C\_05 & C\_06 & .436 & Inequality-reduction argument versus small-business survival concern. \\
Attack & C\_05 & C\_10 & .547 & Inequality-reduction argument versus systemic-critique rejection of the reform. \\
Attack & C\_08 & C\_10 & .350 & Basic-income transition argument versus systemic-critique rejection of the reform. \\
Support & C\_01 & C\_02 & .645 & Two pro-policy worker-protection arguments cohere empirically. \\
Support & C\_01 & C\_04 & .546 & Living-wage/subsidy argument aligns with older-worker income argument. \\
Support & C\_01 & C\_05 & .588 & Living-wage/subsidy argument aligns with inequality-reduction argument. \\
Support & C\_01 & C\_08 & .517 & Living-wage/subsidy argument aligns with basic-income transition argument. \\
Support & C\_02 & C\_04 & .614 & Worker-protection argument aligns with older-worker income argument. \\
Support & C\_02 & C\_05 & .526 & Worker-protection argument aligns with inequality-reduction argument. \\
Support & C\_02 & C\_08 & .416 & Worker-protection argument aligns with basic-income transition argument. \\
Support & C\_03 & C\_06 & .526 & Inflation/fairness concern aligns with small-business survival concern. \\
Support & C\_03 & C\_07 & .364 & Inflation/fairness concern aligns with automation concern. \\
Support & C\_03 & C\_10 & .534 & Inflation/fairness concern aligns with systemic-critique rejection of the reform. \\
Support & C\_04 & C\_05 & .620 & Older-worker income argument aligns with inequality-reduction argument. \\
Support & C\_04 & C\_08 & .527 & Older-worker income argument aligns with basic-income transition argument. \\
Support & C\_05 & C\_08 & .643 & Inequality-reduction argument aligns with basic-income transition argument. \\
Support & C\_06 & C\_07 & .411 & Small-business survival concern aligns with automation concern. \\
Support & C\_06 & C\_09 & .425 & Small-business survival concern aligns with service-redesign concern. \\
Support & C\_06 & C\_10 & .594 & Small-business survival concern aligns with systemic-critique rejection of the reform. \\
Support & C\_07 & C\_09 & .343 & Automation concern aligns with service-redesign concern. \\
Support & C\_07 & C\_10 & .475 & Automation concern aligns with systemic-critique rejection of the reform. \\
Support & C\_09 & C\_10 & .336 & Service-redesign concern aligns with systemic-critique rejection of the reform. \\
\end{longtable}
\endgroup

\subsection{Persuasiveness and Cross-Over Endorsement}

Each consideration also receives a persuasiveness parameter. We considered
simple endorsement rates and item-response discrimination parameters, but both
would make the parameter partly a measure of how already-convinced participants
responded. The final scenario instead uses cross-over endorsement: the share of
participants on the opposing side of a one-dimensional two-parameter IRT
opinion scale who nevertheless agreed with the statement. For pro-policy
considerations, the opposing side is $\theta<0$; for con-policy considerations,
the opposing side is $\theta>0$. For consideration $c$,
\[
\mathrm{cross}_c =
P(\mathrm{agree\ on\ } c \mid \theta \mathrm{\ on\ opposing\ side}).
\]
The observed cross-over rates are then min-max rescaled from 0.3 to 0.9 and
stored as \texttt{persuasiveness} in the simulation scenario.
This makes persuasiveness a measure of empirical reach into the other side of
the opinion space, rather than a measure of raw popularity.

\subsection{Ising Profile Construction}

The agent initialization uses an Ising-style coupling model
\citep{ising1925}. The purpose is to convert sparse participant votes into
complete, empirically coherent consideration profiles for simulation agents.
The selected considerations define a 10-item vote vector. Observed agree and
disagree votes are clamped at $+1$ and $-1$. Passes and missing votes are
imputed using the co-vote couplings among considerations.

For each selected consideration $i$, the external field is derived from the
mean directional vote among participants:
\[
h_i = \operatorname{atanh}(\bar v_i),
\]
with the empirical mean clipped away from $-1$ and $1$. For participant $p$,
let $O_p$ be the set of selected considerations on which the participant cast a
directional vote. Missing or passed consideration $i$ is imputed by a
mean-field expectation:
\[
m_{pi} = \tanh\left(h_i + \sum_{j \in O_p} Q_{ij}v_{pj}\right).
\]
The resulting profile weight $w_{pi}$ is the observed vote when available and
$m_{pi}$ otherwise. Participants are included in the empirical initialization
pool if they have at least three directional votes among the ten selected
considerations. This criterion yields 108 engaged profiles out of 339
participants. Within this engaged subset, participants contribute on average
5.40 observed votes over the selected 10-item set (median $=5$, range $=3$ to
$10$), so 54.0\% of participant-item cells are directly observed and 46.0\%
are imputed by the mean-field step.

Two summary quantities are derived from these profiles. The policy score is
the average direction-weighted profile,
\[
r_p = \frac{1}{10}\sum_i d_i w_{pi},
\]
where $d_i \in \{-1,+1\}$ is the content-coded policy direction of
consideration $i$. The coherence score is based on the pairwise Ising energy,
\[
E_p = -\sum_{i<j} Q_{ij}w_{pi}w_{pj}.
\]
We standardize negative energy over the engaged empirical profile pool, so
higher coherence means a participant's profile aligns more strongly with the
observed co-vote coupling structure.

The resulting engaged profile pool is substantively skewed but not one-sided:
69 participants (63.9\%) are pro-leaning on the direction-weighted policy
score, 31 (28.7\%) are con-leaning, and 8 (7.4\%) fall in an ambivalent band
around zero. On coherence, 22 participants (20.4\%) are in the coherent tail,
59 (54.6\%) in the mixed middle, and 27 (25.0\%) in the ambivalent tail. These
summary statistics justify the paper's use of both policy position and
within-profile coherence as independent dimensions of agent composition.

\subsection{Why These Statements and Not Others}

The selected statements are best understood as cluster representatives, not as
the ten most important statements in the original conversation. Selection is
constrained by three rules: a statement must have enough directional votes to
support pairwise estimation; it must represent one of the semantic clusters;
and, within that cluster, it must have the highest separation power. Omitted
eligible statements are therefore omitted because their semantic cluster is
already represented by a statement with stronger cross-cluster separation in
the co-vote structure. They are not treated as false, irrelevant, or
substantively unimportant.

\begingroup
\scriptsize
\setlength{\tabcolsep}{2pt}
\begin{longtable}{p{0.38in}p{0.42in}p{1.40in}p{3.80in}}
\caption{Moderated statements not selected for the final consideration pool.}\label{tab:omitted-statements}\\
\toprule
Polis & Cluster & Reason & Statement identifier \\
\midrule
\endfirsthead
\toprule
Polis & Cluster & Reason & Statement identifier \\
\midrule
\endhead
\bottomrule
\endfoot
43 & -- & Below threshold: 33 directional votes, one short of eligibility. & It's possible benefits may be cut in exchange for higher wages, along with other benefits. Worth it? \\
1 & 4 & Lower separation than C\_04. & This is a game changer; unintended consequences are too complex to fully understand. \\
2 & 1 & Lower separation than C\_01. & Customers across the board are going to have to pay more. \\
3 & 1 & Lower separation than C\_01. & Critique of the assumption that businesses only have owners' profits to cut. \\
4 & 1 & Lower separation than C\_01. & Anti-corporate-greed argument focused on larger companies. \\
5 & 7 & Lower separation than C\_07. & This will lead to robots. \\
6 & 1 & Lower separation than C\_01. & Smaller businesses may struggle, though the policy is not necessarily bad. \\
7 & 5 & Lower separation than C\_05. & Better implementation alternatives, including subsidies. \\
9 & 4 & Lower separation than C\_04. & Paying \$15 per hour could create a significant boost to the economy. \\
10 & 1 & Lower separation than C\_01. & Companies may not raise prices; wealthy owners can absorb the cost. \\
12 & 4 & Lower separation than C\_04 by a narrow margin. & It is called a living wage for a reason. \\
18 & 6 & Lower separation than C\_06. & Impact on small and medium-sized businesses. \\
24 & 5 & Lower separation than C\_05. & Good move; will boost economy here. \\
25 & 6 & Lower separation than C\_06. & The wage is high enough to cost some needy people their jobs. \\
26 & 5 & Lower separation than C\_05. & Supportive statement suggesting the wage should rise to \$20. \\
28 & 6 & Lower separation than C\_06 by a narrow margin. & This will cause small businesses to go out of business. \\
29 & 5 & Lower separation than C\_05. & Small business will need to make poverty fashionable. \\
34 & 6 & Lower separation than C\_06. & Guaranteed minimum income; remove specific programs. \\
44 & 7 & Lower separation than C\_07. & Minimum wage will not just lead to robots; unionized robots are raised. \\
51 & 4 & Lower separation than C\_04. & The minimum wage should be \$25 per hour. \\
\end{longtable}
\endgroup

Two consequences follow from this selection logic. First, the final
consideration pool should be interpreted as a deliberately compressed map of
argumentative positions. It is designed for experimental control in simulation,
not for exhaustive qualitative representation. Second, several excluded
statements were close alternatives to the selected representative in their
cluster. For example, statement 12 was almost tied with statement 8 in the
older-worker/living-wage cluster, and statement 28 was almost tied with
statement 46 in the small-business cluster. The final pool is therefore
reproducible and empirically grounded, but it should not be overread as the
only defensible qualitative summary of the Polis discussion.
% --- END MERGED BLOCK: methods appendix (method-appendix-semantic-ising.tex) ---
